%%======================================================================
\section{Conclusion}
\label{sec:conclusion}
%%======================================================================

As of March 2026, no bipedal humanoid robot holds FDA clearance or CE marking for
any clinical application. No registered Phase~I, II, or III clinical trial has
enrolled patients for a study involving a full-size bipedal humanoid platform. The
systematic search underlying this review mapped the complete clinical evidence
base for this platform class across 2022--2026, a young body of work still
concentrated on a
single commercial platform (Unitree G1) and, for the clinically most advanced
domain, on a single research cluster (UCSD Yip Laboratory and Johns Hopkins). That
concentration is not evidence of a field in failure; it is evidence of a field in
its earliest phase of organized research. The infrastructure that would enable
distributed clinical investigation (accepted safety standards for bipedal
patient-contact environments, IRB-validated study protocols, and regulatory
pathways with established Special Controls) does not yet exist, and its absence
explains the concentration far better than any deficit of platform capability or
clinical motivation.

This review documents that gap in a form that can guide future work. Prior
reviews of healthcare robotics, including the
closest comparative work in the literature~\cite{cunha2025humanoid}, either
predate the LLM/VLA integration wave that defines the current period, or treat
humanoid and non-humanoid platforms interchangeably,
obscuring the platform-specific barriers and advantages that define the clinical
deployment problem. To our knowledge, this is the first review to integrate the
2022--2026 humanoid-specific literature with a commercial and academic platform
capability assessment, a regulatory pathway analysis, and a per-domain TRL
readiness evaluation grounded in clinical deployment criteria.

Readiness is not uniform across applications. Three roles have a plausible path to
first human-subject validation (TRL~5) within five years: a telemedicine and
telepresence clinical assistant operating in fully teleoperated mode, which
minimizes autonomous patient-contact risk and maps directly onto the established
regulatory logic of physician-controlled surgical robotics; a rehabilitation
movement demonstrator that exploits the mirror neuron mechanism established by
Nguyen et al.~\cite{nguyen2024eeg} and the clinical acceptability demonstrated by
Sobrepera et al.~\cite{sobrepera2022feasibility} without requiring any patient
contact, rendering the bipedal gait safety gap irrelevant to this role; and a
structured eldercare companion and physical assistant following the hybrid
autonomous/teleoperated architecture demonstrated by RHP
Friends~\cite{benallegue2025rhp} and the deployment methodology established by
the SPRING project~\cite{alamedapineda2024spring}. These roles are the most
near-term candidates because they align the strongest available evidence with
tractable technical barriers and comparatively low regulatory risk.

The HCDR framework introduced in Section~\ref{sec:hcdr} makes this
role-differentiated picture precise. By classifying roles along three axes (autonomy
level, patient-contact intensity, and regulatory complexity), the framework
shows that each application is delayed by a different constraint. For the
telemedicine assistant role, the main obstacle is a missing gait safety standard;
for the rehabilitation demonstrator, it is clinical evidence; for eldercare
physical assistance, both are required. This distinction has a practical
consequence: work on gait safety standards and a rehabilitation feasibility trial
can proceed in parallel. The field does not face a single bottleneck, but several
separable bottlenecks that can be addressed through targeted studies.

Two gaps are more consequential than the others identified in
Section~\ref{sec:challenges}. The first is the absence of any accepted performance
standard for bipedal humanoid robots in patient-proximate environments, specifically
for the four categories of gait-induced contact events that ISO/TS~15066 was not
designed to address. Closing this gap does not require a new robot architecture;
it requires biomechanical characterization and a standards-development process.
Without such evidence, IRBs will have little basis for approving physical
patient-contact studies, and regulators will have little basis for defining
Special Controls. This is the most broadly blocking gap in the field. The second
is the absence of registered
human-subject feasibility studies. Cho et al.~\cite{cho2026humanoid} advance the
state of the art by completing a cadaveric surgical study under a defined IRB
protocol, the step immediately preceding first-in-human research. UC San Diego
and Johns Hopkins, which together hold all three papers in the clinical procedures
domain, are well positioned to lead this next step. Beyond the gait safety gap,
no obvious technical reason prevents such a study from being designed in the near
term.

The question confronting the field is no longer whether humanoid robots can
perform clinically meaningful tasks in controlled settings. The evidence reviewed
here indicates that they can. What remains uncertain is whether the research
community, device manufacturers, clinical institutions, and regulatory agencies
will build the sequence of evidence needed to move from laboratory performance to
clinical evidence and, eventually, regulatory precedent. That sequence (safety
characterization, IRB feasibility study, De~Novo submission, and predicate device
establishment) is familiar in medical device development. Applying it to a new
platform class with bipedal dynamics and integrated AI systems will require a
level of institutional coordination that has not yet emerged. Building that
infrastructure is the central task for the next five years.

